% ---------------------------------------------------------------
\section{Introduction}
\label{sec:intro}

\IEEEPARstart{M}{achine-learning} models are trained on historical data and then used
on data that keep changing. When the joint distribution $P(X,Y)$ moves away from the
training distribution, accuracy can fall without any visible error~\cite{gama2014survey,lu2019learning,webb2016characterizing}.
Practitioners usually respond by monitoring the model's \emph{inputs}: a
two-sample test compares a recent window with a reference window, and an alarm is
raised when the difference is significant~\cite{rabanser2019failing,gretton2012kernel,lukats2025benchmark}.
Monitoring inputs is attractive because labels are expensive and often arrive days or
weeks after the prediction is made~\cite{grzenda2020delayed,zliobaite2010change}.

Input monitoring answers a narrower question than the one operators need. A shift in
an irrelevant feature or a seasonal fluctuation can trigger an alarm while accuracy is
unchanged. A change confined to the labelling rule $P(Y\mid X)$ can destroy accuracy
without any change in the inputs. In that case accuracy is not identifiable from
unlabelled data alone~\cite{bayram2022concept,lipton2018detecting}. Confidence-based
estimators of accuracy~\cite{garg2022atc,guillory2021predicting,kivimaki2025confidence}
and drift explainers~\cite{hinder2024surveyB,szucs2025dde} add useful evidence, but
each is usually evaluated in isolation. Visual-analytics systems for drift~\cite{yang2020driftvis,wang2020conceptexplorer}
display that evidence but are rarely scored against verified harm. Their displays also
seldom show which outcomes were still unknown when a decision was made.

\textbf{Research gap.} No study known to the authors evaluates, \emph{in one
label-safe pipeline}, how input alarms, feature-level evidence and calibrated accuracy
estimates relate to verified model harm while labels arrive with an explicit delay.
None also reports where this combination fails. The gap is empirical as well as
architectural. Most evaluations let the monitor (or the evaluator) see labels
earlier than a deployed system could, and they report detection rates without
reporting the harm that went unflagged.

This paper closes that gap with \idm{} (Impact-aware Drift Monitoring with Visual
explanation) and an evaluation protocol whose settings were locked before each run. The contributions are:
\begin{itemize}
  \item A label-safe architecture that separates label-free input evidence from
  calibrated impact evidence and releases ground truth only through a delay-enforcing
  broker (Section~\ref{sec:method}).
  \item A controlled benchmark with known change type and known harm, covering two
  learners, three delays and a 1,500-seed replication (Sections~\ref{sec:setup}--\ref{sec:results}).
  \item Quantified failure modes: under conditional-only change none of 1,205 harmful
  batches is flagged, and closing the dependence blind spot nearly triples false
  alarms (Sections~\ref{sec:results} and~\ref{sec:ablation}).
  \item Evidence from three natural streams and a prospective harm forecaster on where
  delayed labels help and where they do not (Section~\ref{sec:results}).
  \item An aligned evidence view and a fully traceable, rerun-verified artefact
  (Sections~\ref{sec:method} and~\ref{sec:decl}).
\end{itemize}

The remainder of the paper is organised as follows.
Section~\ref{sec:related} reviews related work by theme. Section~\ref{sec:method}
formulates the problem and describes \idm{}. Section~\ref{sec:setup} gives the
experimental setup, and Section~\ref{sec:results} the results. Section~\ref{sec:comparison}
compares the work with its base papers, and Section~\ref{sec:ablation} reports
ablations. Section~\ref{sec:discussion} discusses implications and limitations, and
Section~\ref{sec:conclusion} concludes.\begin{itemize}
  \item A label-safe architecture that separates label-free input evidence from
  calibrated impact evidence and releases ground truth only through a delay-enforcing
  broker (Section~\ref{sec:method}).
  \item A controlled benchmark with known change type and known harm, covering two
  learners, three delays and a 1,500-seed replication (Sections~\ref{sec:setup}--\ref{sec:results}).
  \item Quantified failure modes: under conditional-only change, none of 1,205 harmful
  batches is flagged, and closing the dependence blind spot nearly triples false
  alarms (Sections~\ref{sec:results},~\ref{sec:ablation}).
  \item Evidence from three natural streams and a prospective harm forecaster, which
  shows where delayed labels help and where they do not (Section~\ref{sec:results}).
  \item An aligned evidence view and a fully traceable, rerun-verified artefact
  (Sections~\ref{sec:method},~\ref{sec:decl}).
\end{itemize}

The remainder of the paper is organised as follows.
Section~\ref{sec:related} reviews related work by theme. Section~\ref{sec:method}
formulates the problem and describes \idm{}. Section~\ref{sec:setup} gives the
experimental setup, and Section~\ref{sec:results} the results. Section~\ref{sec:comparison}
compares the work with its base papers, and Section~\ref{sec:ablation} reports
ablations. Section~\ref{sec:discussion} discusses implications and limitations, and
Section~\ref{sec:conclusion} concludes.
